\subsubsection{点双 / 边双 (Biconnected Components)}
同时求桥、边双、点双及点双边集。顶点从 1 编号，原始无向边从 0 编号，以边编号跳过父边，正确支持重边。自环不影响桥／点分割，点双边集约定忽略自环；若题目要求给自环分组，需另行定义。每组 init(n) 后加边、run()；run 可重复运行，时间空间 $O(n+m)$。
\begin{lstlisting}
struct Biconnected {
    int n,timer=0;
    vector<pair<int,int>> edges;
    vector<vector<pair<int,int>>> g;
    vector<int> dfn,low,stack,component,seen;
    vector<char> bridge;
    vector<vector<int>> blocks,block_edges;
    void init(int size){n=size;edges.clear();g.assign(n+1,{});}
    int add(int u,int v){
        int id=edges.size();edges.push_back({u,v});
        if(u!=v){g[u].push_back({v,id});g[v].push_back({u,id});}
        return id;
    }
    void dfs(int u,int parent_edge){
        dfn[u]=low[u]=++timer;
        for(auto [v,id]:g[u]){
            if(id==parent_edge)continue;
            if(!dfn[v]){
                stack.push_back(id);dfs(v,id);low[u]=min(low[u],low[v]);
                if(low[v]>dfn[u])bridge[id]=true;
                if(low[v]>=dfn[u]){
                    blocks.push_back({});block_edges.push_back({});
                    int stamp=blocks.size();
                    for(;;){
                        int e=stack.back();stack.pop_back();
                        block_edges.back().push_back(e);
                        auto [a,b]=edges[e];
                        for(int x:{a,b})if(seen[x]!=stamp){
                            seen[x]=stamp;blocks.back().push_back(x);
                        }
                        if(e==id)break;
                    }
                }
            }else if(dfn[v]<dfn[u]){
                stack.push_back(id);low[u]=min(low[u],dfn[v]);
            }
        }
    }
    void run(){
        timer=0;dfn.assign(n+1,0);low.assign(n+1,0);seen.assign(n+1,0);
        bridge.assign(edges.size(),false);stack.clear();blocks.clear();block_edges.clear();
        for(int u=1;u<=n;u++)if(!dfn[u]){
            if(g[u].empty())blocks.push_back({u}),block_edges.push_back({});
            dfs(u,-1);
        }
        component.assign(n+1,0);int count=0;
        for(int u=1;u<=n;u++)if(!component[u]){
            vector<int> q{u};component[u]=++count;
            while(!q.empty()){
                int x=q.back();q.pop_back();
                for(auto [v,id]:g[x])if(!bridge[id]&&!component[v]){
                    component[v]=count;q.push_back(v);
                }
            }
        }
    }
};
\end{lstlisting}

\paragraph{使用条件与多测} blocks[k] 与 block\_edges[k] 一一对应；只遍历该点双的边集才能维持线性总复杂度，不要反复扫描割点原邻接表。dfs 为递归实现，最大深度 $n$；大规模深链须使用显式栈或足够栈空间。
